\documentclass[12pt,a4paper]{article}
\usepackage[utf8]{inputenc}
\usepackage[T1]{fontenc}
\usepackage[french]{babel}
\usepackage{geometry}
\usepackage{hyperref}
\usepackage{listings}
\usepackage{xcolor}

\geometry{hmargin=2.5cm,vmargin=2.5cm}

% Configuration de la coloration syntaxique pour Java
\lstset{
    language=Java,
    basicstyle=\ttfamily\small,
    keywordstyle=\color{blue}\bfseries,
    stringstyle=\color{red},
    commentstyle=\color{green!60!black},
    numbers=left,
    numberstyle=\tiny\color{gray},
    stepnumber=1,
    numbersep=10pt,
    backgroundcolor=\color{gray!10},
    showspaces=false,
    showstringspaces=false,
    showtabs=false,
    frame=single,
    tabsize=4,
    captionpos=b,
    breaklines=true,
    breakatwhitespace=false,
    title=\lstname,
}

\title{Analyse Approfondie de la Complétion du Compilateur Mini-C \\ \large Projet Sémantique et Traduction des Langages}
\author{Antigravity}
\date{\today}

\begin{document}
\maketitle

\tableofcontents
\newpage

\section*{Introduction}
Ce document présente une analyse extrêmement détaillée et illustrée de l'évolution du dépôt Git pour le projet \textbf{Mini-C}. Depuis son état initial jusqu'à sa complétion, plus de 1700 lignes de code ont été rédigées pour traduire le langage source vers le langage d'assemblage de la machine virtuelle à pile \textbf{TAM}.

La réalisation du projet suit rigoureusement quatre phases : l'analyse syntaxique, l'analyse sémantique (résolution et typage), la gestion de la mémoire, et enfin la génération de code. Pour chaque phase, ce rapport détaille l'intégration des concepts théoriques avec de véritables extraits du code produit.

\newpage
\section{Phase 1 : Analyse Syntaxique et Construction de l'AST}

La première phase du compilateur prend en charge la lecture de la chaîne de caractères (le code source Mini-C), sa validation vis-à-vis d'une grammaire hors contexte (via ANTLR), et surtout la construction d'un Arbre Syntaxique Abstrait (AST).

Le squelette initial de l'analyseur ANTLR générait bien un arbre de dérivation (\textit{Parse Tree}), mais cet arbre contenait beaucoup trop d'informations inutiles pour la compilation (parenthèses, mots-clés, espaces). Le travail a consisté à implémenter le patron de conception \textit{Visitor} via la classe \texttt{ASTBuilder} pour transformer cet arbre de dérivation en un AST pur, composé d'objets métiers (\texttt{Expression}, \texttt{Instruction}, \texttt{Type}).

\subsection{Concept : Les Grammaires Attribuées (TD-03)}
Dans le paradigme des grammaires attribuées, chaque règle de grammaire synthétise un attribut (ici, un sous-arbre de l'AST). Par exemple, lorsque le parseur rencontre une déclaration d'allocation de pointeur (ex: \texttt{malloc(int)}), l'outil ANTLR déclenche une méthode événementielle dans \texttt{ASTBuilder}.

Les ajouts majeurs ont consisté à supporter les types de données complexes, dont les pointeurs. Voici comment le code a été complété dans la classe \texttt{ASTBuilder} :

\begin{lstlisting}[caption={Extrait de ASTBuilder.java - Allocation de pointeurs}]
@Override
public void exitExpressionPointerAllocation(ExpressionPointerAllocationContext ctx) {
    // ctx.type() renvoie le sous-arbre de derivation du type.
    // L'attribut 'unType' a ete synthetise par la visite precedente du nœud Type.
    ctx.uneExpression = new PointerAllocation(ctx.type().unType);
}

@Override
public void exitExpressionPointerAccess(ExpressionPointerAccessContext ctx) {
    // Creation d'un nœud d'acces a un pointeur (*p)
    ctx.uneExpression = new PointerAccess(ctx.expression().uneExpression);
}
\end{lstlisting}

\subsection{Explication de l'architecture}
Dans ce code, l'attribut \texttt{uneExpression} agit comme un \textbf{attribut synthétisé} classique en théorie des langages. Le nœud \texttt{PointerAllocation} va stocker son type enfant (le type de l'objet pointé). Ainsi, lors de l'exécution complète du visiteur, la racine du programme possédera une instance complète de la classe \texttt{Program}, contenant récursivement toutes les listes d'instructions et d'expressions. 

Cette séparation entre la grammaire (ANTLR) et l'AST permet d'isoler la suite du compilateur de toute contrainte textuelle. Le compilateur manipulera uniquement des objets conceptuels, facilitant grandement la sémantique et la génération de code.

\newpage
\section{Phase 2 : Sémantique Statique (Résolution et Typage)}

Un programme syntaxiquement correct n'est pas nécessairement valide. Il peut appeler des fonctions non déclarées, ou additionner un entier avec un booléen. La phase 2 doit parcourir l'AST pour détecter ces anomalies avant d'autoriser la génération du code. C'est l'objectif des méthodes \texttt{resolve()} et \texttt{getType()}.

\subsection{Concept : Résolution et Table des Symboles (TD-04)}
L'AST est initialement un simple arbre. Lorsqu'on utilise une variable \texttt{x}, l'arbre possède un nœud \texttt{IdentifierAccess("x")}. L'étape de \textbf{résolution} consiste à transformer l'arbre en un Graphe Orienté Acyclique (DAG) en reliant l'utilisation de \texttt{x} au nœud de sa déclaration (\texttt{VariableDeclaration("x")}).

Pour ce faire, nous avons implémenté des tables de symboles imbriquées (les environnements lexicaux). L'ajout crucial s'est effectué dans les déclarations de fonctions pour gérer la portée locale des paramètres :

\begin{lstlisting}[caption={Extrait de FunctionDeclaration.java - Résolution}]
@Override
public boolean collectAndPartialResolve(HierarchicalScope<Declaration> _scope) {
    if (_scope.accepts(this)) {
        _scope.register(this); // Enregistre la fonction dans la portee globale
        
        // Creation d'une nouvelle portee locale pour les parametres
        HierarchicalScope<Declaration> paramScope = new SymbolTable(_scope);
        boolean ok = true;
        
        // Resolution des parametres
        for (ParameterDeclaration p : this.parameters) {
            if (paramScope.accepts(p)) {
                paramScope.register(p);
            } else {
                Logger.error("Parameter : " + p.getName() + " is already defined.");
                ok = false;
            }
        }
        // Le corps de la fonction herite de la portee des parametres
        ok &= this.body.collectAndPartialResolve(paramScope, this);
        return ok;
    } else {
        Logger.error("Function : " + this.name + " is already defined.");
        return false;
    }
}
\end{lstlisting}

\subsection{Explication : La Portée Lexicale et le Typage}
Ce bloc de code montre une application parfaite des portées. Un dictionnaire \texttt{SymbolTable} est chaîné à son parent (\texttt{\_scope}). Si une variable n'est pas trouvée localement, la recherche remonte. Si la fonction trouve des paramètres dupliqués (ex: \texttt{void foo(int a, int a)}), l'erreur est immédiatement remontée grâce au booléen \texttt{ok = false}.

Par la suite, la méthode \texttt{getType()} parcourt cet AST résolu en inférant les types de manière ascendante. Les méthodes ajoutées comparent les types de retour, les signatures des appels, garantissant un typage strict, interdisant de ce fait les fameuses \texttt{SemanticsUndefinedException} initialement présentes.

\newpage
\section{Phase 3 : Gestion Mémoire et Allocation}

Une fois l'AST valide (typé et résolu), le compilateur doit anticiper l'exécution sur la machine virtuelle TAM en dimensionnant la mémoire nécessaire pour chaque variable et en leur attribuant une adresse unique (offset par rapport à un registre de base).

\subsection{Concept : Registres de Base et Offsets (TD-06 \& TD-09)}
La TAM gère sa mémoire avec une pile d'exécution. Nous utilisons :
\begin{itemize}
    \item \textbf{SB} (\textit{Stack Base}) pour les variables globales (programme principal).
    \item \textbf{LB} (\textit{Local Base}) pour pointer le début de l'espace local d'un appel de fonction.
\end{itemize}

L'implémentation de \texttt{allocateMemory()} est primordiale. Les paramètres doivent se situer *sous* le registre LB (offsets négatifs), tandis que les variables locales se situent *au-dessus* (offsets positifs, en évitant les 3 mots système réservés à l'appel).

\begin{lstlisting}[caption={Extrait de FunctionDeclaration.java - Allocation mémoire}]
@Override
public int allocateMemory(Register _register, int _offset) {
    // 1. Calculer la taille totale des parametres
    int totalParamSize = 0;
    for (ParameterDeclaration p : this.parameters) {
        totalParamSize += p.getType().length();
    }
    
    // 2. Assigner les offsets negatifs aux parametres (relatifs a LB)
    // Les parametres sont empiles de gauche a droite par l'appelant
    int paramOffset = -totalParamSize;
    for (ParameterDeclaration p : this.parameters) {
        p.setOffset(paramOffset);
        paramOffset += p.getType().length();
    }
    
    // 3. Les variables locales de la fonction (le corps) debutent a LB + 3
    // (Apres l'entête d'appel de 3 mots : Lien Statique, Retour, Ancien LB)
    this.body.allocateMemory(Register.LB, 3);
    
    // Une declaration de fonction ne consomme pas d'espace dans le bloc appelant
    return 0;
}
\end{lstlisting}

\subsection{Explication : La Convention d'Appel}
Cette méthode \texttt{allocateMemory} illustre le cœur du paradigme d'allocation en pile. Le compilateur doit simuler la structure de la TAM :
\begin{itemize}
    \item L'appelant empile ses arguments (taille $T_{param}$).
    \item Le \texttt{CALL} empile 3 mots systèmes (sauvegarde du contexte) et déplace le registre LB juste au-dessus.
    \item L'appelé accède à ses paramètres avec un déplacement vers le bas : on part de \texttt{-totalParamSize}.
    \item L'appelé place ses variables locales après l'entête : l'appel de \texttt{body.allocateMemory(Register.LB, 3)} est l'illustration technique de cette architecture matérielle.
\end{itemize}

\newpage
\section{Phase 4 : Génération de Code Assembleur (Machine Virtuelle TAM)}

La toute dernière étape est la création effective des instructions machines virtuelles. La méthode \texttt{getCode(TAMFactory)} génère récursivement le code de tous les enfants de l'AST et les concatène dans un \texttt{Fragment}.

\subsection{Concept : Fragments, Sauts et Opérations TAM (TD-07)}
Le remplacement des \texttt{SemanticsUndefinedException} par de la véritable logique de compilation se termine ici. Les variables utilisent les offsets calculés en phase 3 pour émettre des \texttt{LOAD (taille) offset[registre]}. 

Le code le plus significatif ajouté concerne l'enrobage des fonctions (déclaration et saut par-dessus) et le nettoyage final de la pile lors du retour.

\begin{lstlisting}[caption={Extrait de FunctionDeclaration.java - Génération de Code}]
@Override
public Fragment getCode(TAMFactory _factory) {
    Fragment _result = _factory.createFragment();
    
    // 1. Saut inconditionnel pour eviter d'executer la fonction lors de sa definition
    int labelNum = _factory.createLabelNumber();
    String endLabel = "end_func_" + this.name + "_" + labelNum;
    _result.add(_factory.createJump(endLabel));
    
    // 2. Generer le corps de la fonction (sans POP terminal, le RETURN gere ça)
    Fragment _body = this.body.getCodeWithoutCleanup(_factory);
    
    int totalParamSize = 0;
    for (ParameterDeclaration p : this.parameters) {
        totalParamSize += p.getType().length();
    }
    
    // 3. Generer l'instruction RETURN (tailleResultat) (tailleParametres)
    int returnSize = this.type.length();
    _body.add(_factory.createReturn(returnSize, totalParamSize));
    
    // 4. Etiquetage du point d'entree
    _body.addPrefix(this.name);
    _result.append(_body);
    
    // 5. Etiquetage de sortie pour le saut inconditionnel
    _result.addSuffix(endLabel);
    return _result;
}
\end{lstlisting}

\subsection{Explication : Automatisation Assembleur}
Dans cette phase, on instruit directement le processeur TAM. 
\begin{itemize}
    \item Le \texttt{JUMP(endLabel)} protège la fonction : le flux principal d'exécution d'un programme Mini-C ne doit pas glisser accidentellement dans le code de la fonction. 
    \item L'étiquette (prefix) assigne le nom textuel de la fonction au point d'entrée, ce qui permet au \texttt{FunctionCall.getCode()} (non illustré ici) d'utiliser simplement l'instruction assembleur \texttt{CALL (nom\_fonction)}.
    \item L'instruction \texttt{RETURN} joue un rôle crucial : non seulement elle redonne le contrôle à l'appelant, mais elle détruit le cadre de la pile locale, et \textbf{surtout}, dépile les paramètres fournis, protégeant ainsi le programme appelant d'une fuite mémoire.
\end{itemize}

Ce patron garantit que tout appel à une fonction, qu'elle soit récursive ou imbriquée, nettoie proprement son environnement. La totalité du \textit{back-end} est ainsi rendue fonctionnelle.

\section*{Conclusion Générale}
La mutation du projet vierge au compilateur final témoigne d'une approche \textit{Top-Down}. L'abstraction (les arbres) est progressivement concrétisée : l'arbre est typé et lié (Phase 2), géométriquement dimensionné dans la mémoire (Phase 3), puis mathématiquement traduit en actions processeur (Phase 4). L'ensemble du système résout élégamment le défi du "dangling pointer", de l'allocation sur la pile, et de la validation des types pour le Mini-C.

\end{document}
